\clearpage
\begin{trapbox}
\textbf{C087 八个高频陷阱与最小反例}

\textbf{1. 把竖直焦点弦漏掉。}
$y=k(x-c)$ 只表示有限斜率。$x=c$ 的斜率不存在，却合法且恰是同支最短弦。
\textbf{反例：}$a=3,b=4,c=5$，竖直弦长 $32/3$。

\textbf{2. 以为 $k=\pm b/a$ 是切线。}
\begin{center}% a=3,b=4,c=5,k=4/3; sole finite P=(17/5,-32/15).
\begin{tikzpicture}[font=\small,>=stealth]
\draw[gray!45,rounded corners=2pt] (0,0) rectangle (14.2,9.5);
\node at (7.1,8.85) {边界不是切线：过焦点的直线平行渐近线，只有一个有限交点};
\begin{scope}[shift={(7.1,4.7)},scale=.48]
\draw[gray,->] (-7,0)--(7.8,0) node[right] {$x$};
\draw[gray,->] (0,-6.4)--(0,6.5) node[above] {$y$};
\draw[blue!75!black,thick,domain=-6:6,samples=100,variable=\v] plot ({3*sqrt(1+\v*\v/16)},\v);
\draw[blue!75!black,thick,domain=-6:6,samples=100,variable=\v] plot ({-3*sqrt(1+\v*\v/16)},\v);
\draw[gray,dashed,thick] (-4.5,-6)--(4.5,6);
\draw[red!75!black,very thick] (.5,-6)--(7.3,{4/3*(7.3-5)});
\fill (5,0) circle(2pt) node[below right=6pt] {$F_2$};
\fill (3.4,{-32/15}) circle(2pt) node[left=8pt] {$P$};
\end{scope}
\node[anchor=west,text=gray!80!black] at (.8,7.7) {灰虚线：$y=\frac43x$};
\node[anchor=west,text=red!75!black] at (.8,6.75) {红实线：$y=\frac43(x-5)$};
\node at (7.1,.65) {联立后二次项消失：$160t+256=0$；不是两根相等，而是只剩一次方程。};
\end{tikzpicture}
\end{center}
它们只是\textbf{平行渐近线的过焦点直线}，联立方程二次项消失，
只有一个\textbf{有限交点}；不能把它当作两点重合的切线。
\textbf{反例：}$a=3,b=4,c=5,k=4/3$，联立得 $160t+256=0$，
仅 $t=-8/5$ 一解。

\tcbbreak
\textbf{3. 同支与异支最小值互换。}
\begin{center}% 精确坐标：H: x^2/9-y^2/16=1, c=5, b^2/a=16/3.
% 两图中蓝色曲线由解析式 x=+-3 sqrt(1+y^2/16) 绘制，非手绘。
\begin{tikzpicture}[x=0.39cm,y=0.39cm,>=stealth, font=\small]
\begin{scope}
  \draw[->,gray] (-6.8,0)--(7.3,0) node[right]{$x$};
  \draw[->,gray] (0,-7.8)--(0,7.8) node[above]{$y$};
  \draw[thick,blue!75!black,domain=-7:7,samples=90] plot ({3*sqrt(1+\x*\x/16)},{\x});
  \draw[thick,blue!75!black,domain=-7:7,samples=90] plot ({-3*sqrt(1+\x*\x/16)},{\x});
  \draw[very thick,red!75!black] (5,-5.33333333)--(5,5.33333333);
  \fill (5,0) circle (2.1pt) node[below right]{$F_2$};
  \fill[red!75!black] (5,5.33333333) circle (1.8pt) node[above left=5pt]{$A$};
  \fill[red!75!black] (5,-5.33333333) circle (1.8pt) node[below left=5pt]{$B$};
  \node[align=center] at (0,-9.0) {同支焦点弦\quad $x=c$\\最短 $2b^2/a$};
\end{scope}
\begin{scope}[xshift=6.6cm]
  \draw[->,gray] (-6.8,0)--(7.3,0) node[right]{$x$};
  \draw[->,gray] (0,-7.8)--(0,7.8) node[above]{$y$};
  \draw[thick,blue!75!black,domain=-7:7,samples=90] plot ({3*sqrt(1+\x*\x/16)},{\x});
  \draw[thick,blue!75!black,domain=-7:7,samples=90] plot ({-3*sqrt(1+\x*\x/16)},{\x});
  \draw[very thick,red!75!black] (-3,0)--(3,0);
  \fill (5,0) circle (2.1pt) node[above right]{$F_2$};
  \fill[red!75!black] (-3,0) circle (1.8pt) node[above left=5pt]{$A$};
  \fill[red!75!black] (3,0) circle (1.8pt) node[above right=5pt]{$B$};
  \node[align=center] at (0,-9.0) {异支焦点弦\quad $y=0$\\最短 $2a$};
\end{scope}
\end{tikzpicture}
\end{center}
异支跨越两个顶点之间的实轴间距：下界 $2a$；
同支关于焦点上下分开：下界 $2b^2/a$。两数没有固定大小关系，
例如 $a=3,b=4$ 时 $32/3>6$；若 $a=4,b=2$，则 $2b^2/a=2<8=2a$。

\textbf{4. 没有给分母加绝对值。}
当 $k=0$ 时 $a^2k^2-b^2=-b^2$，没有绝对值会得出负弦长。

\textbf{5. 用 $L\to2b^2/a$ 断言取不到。}
“有限斜率 $|k|\to\infty$”只是逼近竖直情况；真实竖直线\textbf{存在}，
所以同支范围是\textbf{左闭}的 $[2b^2/a,+\infty)$。

\textbf{6. 只看 $k$ 正负来判别分支。}
分类依据 $|k|$ 相对于 $b/a$，不是 $k>0$ 或 $k<0$。

\textbf{7. 混淆焦点弦与过原点的中心弦。}
本结论直线必须经过 $F_1$ 或 $F_2$。一般中心弦不能直接套公式。

\textbf{8. 取值范围中的 $+\infty$ 不是一个可取的弦长。}
当 $k\to\pm b/a$ 时有限弦长无界增长；在临界斜率处没有两个有限交点。

\textbf{最低门槛纠错检验：}把 $k=0$ 带入必须得到 $2a$；
把 $\theta=\pi/2$ 带入必须得到 $2b^2/a$；
把 $k=\pm b/a$ 带入必须被判定为不合法。
\end{trapbox}
